% Generated draft: requires mathematical and source review.
\documentclass[11pt]{article}
\usepackage[margin=25mm]{geometry}
\usepackage{fontspec}
\setmainfont{Latin Modern Roman}
\usepackage{amsmath,amssymb,mathtools,unicode-math}
\setmathfont{Latin Modern Math}
\usepackage{enumitem,needspace,xurl,hyperref}
\setlength{\parindent}{0pt}
\setlength{\parskip}{5pt}
\setlength{\emergencystretch}{3em}
\newcommand{\N}{\mathbb{N}}
\newcommand{\Z}{\mathbb{Z}}
\newcommand{\Q}{\mathbb{Q}}
\newcommand{\R}{\mathbb{R}}
\newcommand{\C}{\mathbb{C}}
\newcommand{\F}{\mathbb{F}}
\newcommand{\A}{\mathbb{A}}
\newcommand{\PP}{\mathbb{P}}
\newcommand{\E}{\mathbb{E}}
\newcommand{\eps}{\varepsilon}
\newcommand{\dd}{\,\mathrm{d}}
\begin{document}
\section*{UnsolvedMath ID: 801138}
\textbf{Source problem code:} catalogue entry
\subsection{Definitions and mathematical statement}
Under the standard tracial interpretation of the notation, let \(N\) be a type \(\mathrm{II}_1\) factor with normalized faithful normal trace \(\tau\), let \(I\) be an index set, and let \(\mathcal U\) be an ultrafilter on \(I\). A type \(\mathrm{II}_1\) factor is an infinite-dimensional finite von Neumann algebra with center \(\C 1\). Write \(\|x\|_{2,\tau}=\tau(x^*x)^{1/2}\). Its tracial ultrapower is
\[
N^{\mathcal U}
=
\ell^\infty(I,N)\big/\mathcal I_{\mathcal U},
\qquad
\mathcal I_{\mathcal U}
=
\left\{(x_i)_{i\in I}:\lim_{i\to\mathcal U}\|x_i\|_{2,\tau}=0\right\},
\]
where \(\ell^\infty(I,N)\) consists of families bounded in operator norm.

For a discrete group \(G\), its group von Neumann algebra is
\[
L(G)=\{\lambda(g):g\in G\}'',
\qquad
\lambda(g)\delta_h=\delta_{gh}
\quad(g,h\in G),
\]
acting on \(\ell^2(G)\), with the double commutant taken in \(B(\ell^2(G))\).

\textbf{Source question.} Can \(N^{\mathcal U}\) ever be isomorphic to a group von Neumann algebra? For specified \(N\) and \(\mathcal U\), the standard discrete-group formulation asks whether
\[
\exists\,G\text{ a discrete group}\quad N^{\mathcal U}\cong L(G),
\]
where \(\cong\) denotes a normal unital \(*\)-isomorphism of von Neumann algebras.

\textbf{Unresolved qualifications.} The supplied record does not define \(N\) or \(\mathcal U\). The type \(\mathrm{II}_1\) and tracial interpretations above are suggested by the source-list title, but their precise hypotheses are not supplied. In particular, the record does not specify separability of \(N\), the index set, whether \(\mathcal U\) must be nonprincipal, or whether \(N\) and \(\mathcal U\) are fixed or existentially quantified by ``ever.'' The discrete-group convention is stated explicitly here; the record does not specify the intended class of groups. Thus the displayed formulation is conditional on these conventions, and the full intended quantifier scope remains unresolved.
\subsection{Short English statement}
The question asks whether an ultrapower of a type \(\mathrm{II}_1\) factor can be isomorphic to a group von Neumann algebra. The supplied record leaves the factor, ultrafilter hypotheses, group class, and scope of ``ever'' unspecified. Its open status is a supplied registry label.
\subsection{Sources}
\begin{itemize}
\item Local UnsolvedMath registry, permanent ID \(801138\), source problem code \texttt{catalogue entry}, title ``Ultraproduct II\_1 Factors --- Problem 39.'' The supplied record attributes the problem to Elayavalli (2025) and labels it open; this status has not been independently verified.
\item Elayavalli, \emph{50 Open Problems: Ultraproduct II\_1 Factors} (2025), Problem 39. Supplied source link: \url{https://arxiv.org/abs/2511.20377}. The supplied provenance records access on 2026-10-16 and extraction by \texttt{arxiv\_latex}; these are supplied provenance claims, not a claim of independent source consultation.
\end{itemize}
\end{document}
